\chapter{Core analysis}\label{chap:core}

This chapter establishes the primary estimate
(Proposition~\ref{prop:primary-box}) by expanding $\log\psi$ to fifth
order, bounding each cumulant correction, and integrating over $\corereg$.

\section{Cumulant identities}

\begin{lemma}[Triangle formula]\label{lem:triangle}
\lean{triangle_formula}\leanok
For every $\lambda\in\R^d$,
$T(\lambda)=\sum_{1\le i<j<k\le n}\lambda_{ij}\lambda_{ik}\lambda_{jk}$.
\end{lemma}

\begin{proof}\leanok
Expand $\E[X_\lambda^3]$ and observe that the only surviving multigraphs
(every vertex with even multiplicity) are triangles. Each triangle
contributes $3!$ ordered triples, giving $\E[X_\lambda^3]=6\sum\lambda_{ij}\lambda_{ik}\lambda_{jk}$.
Since $T=\frac16\E[X_\lambda^3]$, the claim follows.
\end{proof}

\begin{lemma}[Fourth cumulant]\label{lem:ordered-cycle}
\lean{fourth_cumulant_identity}\leanok
\uses{lem:triangle}
$Q(\lambda)
=-\frac{1}{12}\sum_e\lambda_e^4+\frac18\mathcal{C}_4(\lambda)$,
where $\mathcal{C}_4(\lambda)=\sum_{\text{ordered 4-cycles}}\lambda_{i_1 i_2}\lambda_{i_2 i_3}\lambda_{i_3 i_4}\lambda_{i_4 i_1}$.
\end{lemma}

\begin{proof}\leanok
Expand $\E[X_\lambda^4]=\sum_{e_1,\dots,e_4}\lambda_{e_1}\cdots\lambda_{e_4}
\E[\chi_{e_1}\cdots\chi_{e_4}]$.
The expectation vanishes unless every vertex index has even multiplicity.
Three types survive:
(i)~all four edges equal: $\sum_e\lambda_e^4$;
(ii)~two distinct pairs: $6\sum_{e<f}\lambda_e^2\lambda_f^2$;
(iii)~four distinct edges forming a 4-cycle: $3\mathcal{C}_4(\lambda)$.
Thus $\E[X_\lambda^4]=\sum_e\lambda_e^4+6\sum_{e<f}\lambda_e^2\lambda_f^2
+3\mathcal{C}_4$.
Since $3s(\lambda)^2=3\sum_e\lambda_e^4+6\sum_{e<f}\lambda_e^2\lambda_f^2$,
subtracting gives $\kappa_4=-2\sum_e\lambda_e^4+3\mathcal{C}_4$, and
dividing by~$24$ yields the formula for $Q$.
\end{proof}

\begin{lemma}[Peeling identity]\label{lem:peeling}
\lean{quarticCorr_peel_last}\leanok
\uses{lem:ordered-cycle}
Let $A$ be a symmetric zero-diagonal $n\times n$ matrix, $B$ the
$(n{-}1)\times(n{-}1)$ submatrix deleting row/column $n$,
$x\in\R^{n-1}$ the last column, and $M(B)=B^2-\diag(B^2)$. Then
\[
  Q_n(A)=Q_{n-1}(B)+\tfrac12 x^\top M(B)x-\tfrac{1}{12}\sum_{a=1}^{n-1}x_a^4.
\]
\end{lemma}

\begin{proof}\leanok
The diagonal quartic part splits as
$\sum_{i<j\le n}A_{ij}^4=\sum_{a<b\le n-1}B_{ab}^4+\sum_a x_a^4$.
Cycles through vertex $n$ correspond to $x^\top M(B)x$;
cycles avoiding $n$ are $\mathcal{C}_{4,n-1}(B)$.
Combine with the $1/8$ and $-1/12$ coefficients.
\end{proof}

\section{Quartic and quintic bounds}

\begin{proposition}[Quartic bounds]\label{prop:quartic-bounds}
\lean{correctedCore_quarticCore_gap_abs_le_gaussianF_short}\leanok
\uses{lem:peeling,lem:gaussian-quadratic,lem:inner-core,lem:hc}
There exist $c_1,C_1>0$ such that if $d/t\le c_1$ and $nd^2/t^2\le c_1$,
then
\[
  \int_{\corereg}|e^{4tQ(\lambda)}|^2 e^{-2t\|\lambda\|^2}\,d\lambda
  \le C_1\core(d,t),
\]
and
\[
  \int_{\corereg}|e^{4tQ(\lambda)}-1|^2 e^{-2t\|\lambda\|^2}\,d\lambda
  \le C_1\frac{d^2}{t^2}\core(d,t).
\]
\end{proposition}

\begin{proof}\leanok
\uses{lem:peeling,lem:gaussian-quadratic,lem:inner-core,lem:hc}
\emph{First estimate.}
Set $I_n^{(\beta)}=\int_{\corereg}e^{-2t\|\lambda\|^2}e^{\beta tQ}\,d\lambda$
with $\beta=8$ (since $|e^{4tQ}|^2=e^{8tQ}$).
By Lemma~\ref{lem:peeling}, $Q_n(A)=Q_{n-1}(B)+\frac12 x^\top M(B)x
-\frac{1}{12}\sum x_a^4$.
Dropping the negative quartic and fixing $B$, the $x$-fiber is
$\int_{\R^{n-1}}\exp(-2t\|x\|^2+\frac{\beta t}{2}x^\top M(B)x)\,dx$.
On $\corereg$, $\|M(B)\|_{\mathrm{op}}\le 2s(B)\le 2d/t$,
so $I-\frac\beta4 M(B)$ is positive definite for small $d/t$.
Lemma~\ref{lem:gaussian-quadratic} evaluates the Gaussian integral;
the determinant satisfies
$\det(I-\frac\beta4 M(B))^{-1/2}\le\exp(C_\beta d^2/t^2)$
since $\operatorname{tr}(M(B)^2)\le 4s(B)^2$ and $\sum_j u_j=0$.
Peeling off all $n$ vertices gives
$I_n^{(\beta)}\le \exp(C_\beta nd^2/t^2)\,F(d,t)$,
bounded by $C_1\core(d,t)$ when $nd^2/t^2\le c_1$
(using Lemma~\ref{lem:inner-core}).

\emph{Second estimate.}
Use $|e^u-1|^2\le u^2(1+e^{2u})$ with $u=4tQ$.
The term $\int Q^2 e^{-2t\|\lambda\|^2}$ equals
$F(d,t)\E_{\gamma_t}[Q^2]$, where $\E_{\gamma_t}[Q^2]=O(d^2/t^4)$
by expanding $Q=-\frac{1}{12}D+\frac18 C$ and using
$\E[DC]=0$ (parity) and $\E[C^2]=O(d^2/t^4)$ (partner counting).
The cross term $\int Q^2 e^{8tQ}e^{-2t\|\lambda\|^2}$ is bounded
by Cauchy--Schwarz, the first estimate with $\beta=16$, and
hypercontractivity (Lemma~\ref{lem:hc}) applied to the degree-4
polynomial $Q$. Multiplying by $16t^2$ gives the result.
\end{proof}

\begin{lemma}[Quintic bound]\label{lem:quintic}
\lean{quinticP5_pointwise_bound}\leanok
\uses{lem:hc,lem:inner-core}
There is $C_3'>0$ such that
\[
  \int_{\corereg}P(\lambda)^2 e^{-2t\|\lambda\|^2}\,d\lambda
  \le C_3'\frac{n^5}{t^5}\core(d,t).
\]
\end{lemma}

\begin{proof}\leanok
\uses{lem:hc,lem:inner-core}
By definition $\kappa_5=\E[X_\lambda^5]-10\E[X_\lambda^3]\E[X_\lambda^2]$,
so the disconnected 5-edge multigraphs (triangle $\times$ double edge) cancel.
Thus $P$ is supported on \emph{connected} even 5-edge multigraphs.
Such a graph has sum of degrees $10$, minimum degree $2$, hence at
most $5$ vertices: choosing $\le 5$ vertices from $[n]$ gives $O(n^5)$
monomials.

Under $\gamma_t$ (variance $(4t)^{-1}$), each monomial has degree 5.
A pair $(M,M')$ survives $\E_{\gamma_t}[MM']$ only if every edge
has even total degree, i.e., $\mathcal O(M')=\mathcal O(M)$
(same odd-support graph). The odd support is either a 5-cycle
(unique monomial, unique partner) or a triangle on $\{i,j,k\}$
(partner types: $(3,1,1)$ on the triangle---$3$ choices---or $(1,1,1,2)$
with a pendant edge to a new vertex $\ell$---$O(n)$ choices).
This gives $O(n^5)$ surviving pairs.
Each degree-10 moment is $O(t^{-5})$; multiplying by $F(d,t)$
and using Lemma~\ref{lem:inner-core} proves the claim.
\end{proof}

\section{Log-expansion and small-ball decay}

\begin{lemma}[Log-expansion]\label{lem:sixth}
\lean{psi_pow_sub_correctedCoreIntegrand_norm_le_uniform}\leanok
\uses{lem:realpart,lem:hc}
There are $c_2,C_2>0$ such that whenever $s(\lambda)\le c_2$,
\[
  \log\psi(\lambda)
  =-\tfrac12 s(\lambda)-iT(\lambda)+Q(\lambda)+iP(\lambda)+E_6(\lambda)
\]
with $|E_6(\lambda)|\le C_2 s(\lambda)^3$.
\end{lemma}

\begin{proof}\leanok
\uses{lem:realpart,lem:hc}
Set $K(u)=\log\E[e^{iuX_\lambda}]$ and Taylor-expand to fifth order.
Lemma~\ref{lem:realpart} ensures $\RePart\psi(u\lambda)\ge 3/4$,
so the principal branch of $\log$ is well-defined on $[0,1]$.
The sixth derivative $K^{(6)}$ involves products $m(u)^{-k}\prod m^{(j_\ell)}(u)$;
each factor is bounded using $|m(u)|\ge 3/4$ and hypercontractivity
(Lemma~\ref{lem:hc}) applied to the degree-2 polynomial $X_\lambda$.
\end{proof}

\begin{lemma}[Small-ball decay]\label{lem:small-ball}
\lean{small_ball_gaussian_decay_uniform}\leanok
\uses{lem:sixth}
There exists $r_0\in(0,\pi/4)$ such that
$|\psi(\lambda)|^{4t}\le e^{-3t\|\lambda\|^2/2}$
for every $t\ge 1$ and $\lambda\in\Dd_{r_0}$.
\end{lemma}

\begin{proof}\leanok
\uses{lem:sixth}
By the log-expansion (Lemma~\ref{lem:sixth}),
$\RePart\log\psi=-\frac12 s+Q+\RePart E_6$.
Show $Q=O(s^2)$ using $|\mathcal{C}_4|\le \operatorname{tr}(A^2)^2=4s^2$
and $|E_6|\le C_2 s^3$. For small $r_0$, the corrections are absorbed
by the leading $-\frac12 s$ term.
\end{proof}

\section{Cubic phase bound}

\begin{lemma}[Cubic phase bound]\label{lem:cubic}
\lean{cubic_core_second_order_gap_uniform}\leanok
\uses{lem:inner-core,lem:gaussian-radial,lem:triangle}
For all $d,t\ge 1$,
\[
  \RePart\int_{\corereg}e^{-2t\|\lambda\|^2}e^{-4itT(\lambda)}\,d\lambda
  \ge \bigl[1-O(n^3/t)\bigr]\core(d,t).
\]
\end{lemma}

\begin{proof}\leanok
\uses{lem:inner-core,lem:gaussian-radial,lem:triangle}
Since $T$ is odd and $\corereg$ is symmetric, the imaginary part vanishes.
The real part satisfies $1-\cos(4tT)\le 8t^2 T^2$. By the triangle formula
(Lemma~\ref{lem:triangle}), $T^2$ is a degree-6 polynomial. Gaussian
radial moments (Lemma~\ref{lem:gaussian-radial}) bound the integral of
$T^2$ against the Gaussian weight, giving an $O(n^3/t)$ correction.
\end{proof}

\section{Primary estimate}

\begin{proposition}[Primary estimate]\label{prop:primary-box}
\lean{prop_primary_box}\leanok
\uses{lem:sixth,lem:small-ball,lem:cubic,prop:quartic-bounds,lem:quintic,lem:inner-core}
Set $K_n:=2^{2d-n+1}(2\pi)^{-d}$.
For all sufficiently large $n$ and $t\ge C_0 n^3$,
\[
  K_n\RePart\int_{\corereg}\psi(\lambda)^{4t}\,d\lambda
  =\left(1-\frac{\binom{n}{3}}{8t}
  +O\!\left(\frac{n^2}{t}+\frac{n^{5/2}}{t^{3/2}}+\frac{n^6}{t^2}
  +e^{-cd}\right)\right)\Ahat.
\]
\end{proposition}

\begin{proof}\leanok
\uses{lem:sixth,lem:cubic,prop:quartic-bounds,lem:quintic,lem:inner-core}
By the log-expansion (Lemma~\ref{lem:sixth}), on $\corereg$ we write
$\psi(\lambda)^{4t}=e^{-2t\|\lambda\|^2}\,e^{-4itT}\,e^{4tQ}\,e^{4itP}\,e^{4tE_6}$
with $|E_6|\le C_2 s^3$.
Define intermediate integrals:
\[
  M_{\mathrm{core}}:=\int_{\corereg}e^{-2t\|\lambda\|^2}e^{-4itT}e^{4tQ}e^{4itP}\,d\lambda,
  \qquad
  L_{\mathrm{core}}:=\int_{\corereg}e^{-2t\|\lambda\|^2}e^{-4itT}e^{4tQ}\,d\lambda.
\]

\emph{Step 1: remove $E_6$.}
$|\psi^{4t}-M_{\mathrm{core}}\cdot(\text{Gaussian})|\le |e^{4tE_6}-1|\le C t s^3$,
and $\int s^3 e^{-2t\|\lambda\|^2}\le C(d/t)^3 F(d,t)$
by Gaussian radial moments.

\emph{Step 2: remove $P$.}
$|M_{\mathrm{core}}-L_{\mathrm{core}}|\le |e^{4itP}-1|\le 4t|P|$.
By Cauchy--Schwarz and Lemma~\ref{lem:quintic},
$\int |P|\,e^{-2t\|\lambda\|^2}\le C n^{5/2}/t^{5/2}\cdot\core(d,t)^{1/2}$.
This contributes $O(n^{5/2}/t^{3/2})\Ahat$.

\emph{Step 3: remove $Q$.}
$|L_{\mathrm{core}}-\int e^{-2t\|\lambda\|^2}e^{-4itT}|\le
|e^{4tQ}-1|$ integrated against the Gaussian weight.
Proposition~\ref{prop:quartic-bounds} gives
$\int |e^{4tQ}-1|^2 e^{-2t\|\lambda\|^2}\le C d^2/t^2\cdot\core(d,t)$,
contributing $O(d/t)=O(n^2/t)$ by Cauchy--Schwarz.

\emph{Step 4: cubic phase.}
By Lemma~\ref{lem:cubic},
$\RePart\int e^{-2t\|\lambda\|^2}e^{-4itT}\ge [1-Cn^3/t]\core(d,t)$.
The correction $\binom{n}{3}/(8t)$ comes from the exact computation
$\E_{\gamma_t}[T^2]=\binom{n}{3}/(64t^3)$ (via the triangle formula)
and $1-\cos(4tT)\approx 8t^2T^2$.

Combining all four steps and using $\core(d,t)=(1+O(e^{-cd}))F(d,t)$
(Lemma~\ref{lem:inner-core}) yields the claimed expansion.
\end{proof}
